% Keep styling.sty beside this file. Compile with:
% latexmk -xelatex -synctex=1 template.tex (or, 
% for my purposes, compile to preview with VimtexCompile in neovim)
% if youre using VSC with some built in plugin you'll probably need to add 
% {
%   "latex-workshop.latex.recipe.default": "latexmk (xelatex)",
%   "latex-workshop.latex.autoBuild.run": "onSave",
%   "latex-workshop.view.pdf.viewer": "tab",
%   "latex-workshop.view.pdf.invertMode.enabled": "never"
% }
% also it is necessary to download the right fonts. Just look 
% at error messages to find which ones or something

\documentclass[11pt]{article}
\usepackage{styling} % Dark by default, with the clickable THEME switch.
% \usepackage{graphicx} % Uncomment if adding external images.

\title{Title of the Notes}
\author{Remote Group}
\date{For November 15th} % Use \date{} to hide the date.

\begin{document}
\maketitle
% \tableofcontents
% \newpage

\begin{prob*}[Munkres 55.4]
    Suppose you are given the fact that 
    for each $n$, there is no retraction 
    $r : B^{n+1} \to S^n$. Prove the following: 
    \begin{enumerate}[label=(\alph*)]
        \item The identity map $i : S^n \to S^n$ is not
            nullhomotopic. 
        \item The inclusion map $j : S^n \to 
            \mathbb R^{n+1} - 0$ is not nullhomotopic. 
        \item Every nonvanishing vector field 
            on $B^{n+1}$ points directly outward at 
            some point of $S^n$, and directly inward 
            at some point of $S^n$. 
        \item Every continuous map $f : B^{n+1} \to B^{n+1}$ 
            has a fixed point. 
        \item Every $n + 1$ by $n+1$ matrix with 
            positive real entries has a positive 
            eigenvalue. 
        \item If $h : S^n \to S^n$ is nullhomotopic, then 
            $h$ has a fixed point and $h$ maps 
            some point $x$ to its antipode $-x$.
    \end{enumerate}
\end{prob*}
\begin{proof}[Proof of (a)]
    We will prove by contrapositive. 

    Apologies if there are some slips in the rigor here, 
    this was hard. 

    Suppose the identity $i : S^n \to S^n$ 
    is nullhomotopic. 

    Then there is a continuous map 
    \[
        H : S^n \times I \to S^n 
    \]
    such that 
    \[
        H(x,0) = x, \ H(x,1) = p 
    \]
    for some fixed $p \in S^n$. 

    Starting with $S^n \times I$ we can impose 
    $(x,1) \sim (y,1)$ for all 
    $x,y \in S^n$. $H$ is then constant on all 
    equivalence classes, letting us factor 
    \[
        \begin{tikzcd}
            S^n \times I \ar[r, "H"] \ar[dr, "q"] & S^n \\ 
            & (S^n \times I)/(\sim) \ar[u, "\varphi"']
        \end{tikzcd}
    \]
    giving a continuous map $\varphi$ for which the 
    diagram commutes ($q$ is $q(x,t) = [(x,t)]$). 
    
    For this map we have that 
    \[
        \varphi \bigm|_{S^n \times 0} 
        = H\bigm|_{S^n} = \id_S^n.  
    \]

    Now it remains to show that 
    \[
        (S^n \times I)/\sim 
        \ \cong B^{n+1}. 
    \]
    Take $f : (S^n \times I)/\sim \ \longrightarrow 
    B^{n+1}$ by 
    \[
        f([(x_1,x_2,\ldots,x_{n+1}),t]) = 
        (1-t)(x_1,x_2,\ldots,x_{n+1}).  
    \]
    Then the point $p$ corresponding to the 
    nontrivial equivalence class is the origin. 

    The map is a well-defined homeomorphism (I swear!).

    This gives a retraction map from $B^{n+1} \to S^n$ 
    which we were told wasn't possible! 
\end{proof}
\begin{proof}[Proof of (b)]
    I'm sure one could argue this from the fact that 
    $S^n \simeq \mathbb R^{n+1} \setminus 0$, but instead,
    take the retraction 
    \[
        r : \mathbb R^{n+1} \setminus 0 
        \to S^n, \ r(x) = \frac{x}{\|x\|}. 
    \]
    Observe that 
    \[
        r \circ j = \id_{S^n}. 
    \]
    Suppose, for contradiction, that $j$ is nullhomotopic. 
    Then 
    \[
        j \simeq c_p
    \]
    for some constant map $c_p : S^n \to \mathbb R^{n+1} 
    \setminus 0$. 

    Composing with $r$, we get 
    \[
        \id_{S^n} = r \circ j \simeq r \circ c_p, 
    \]
    where $r \circ c_p$ is constant, contradicting (a).    
\end{proof}
\begin{prob*}
    A nonvanishing vector field on 
    $B^{n+1}$ is a continuous map 
    \[
        v : B^{n+1} \to \mathbb R^{n+1}
    \]
    such that 
    \[
        v(x) \neq 0 \ \forall x \in B^{n+1}. 
    \]
\end{prob*}

\end{document}